\section{The observation and the main results}\label{sec:setting}
Near the onset of a prescribed billiard itinerary, phase volume contains
a residual-time factor as well as a collision-section flux. Keeping
intrinsic endpoint positions rather than integrating them out makes
this structure visible. Two density functions recover the local action
and its volume normalization. The normalization has a second use:
after analytic boundary recovery it determines the covolume of the
unknown period lattice. This connects the local observation to a
finite-index global reconstruction without supplied lattice coordinates.

\subsection{Local records}
A periodic dispersing table is the complement of
\[
             \bigcup_{i=1}^{r}\ \bigcup_{\lambda\in\Lambda}
                         (C_i+\lambda)\subset\R^2,
\]
where $\Lambda$ is a full-rank lattice and the $C_i$ are compact strictly
convex bodies with smooth positively curved boundaries. All translates
are disjoint with positive separation. There is one representative per
obstacle orbit, and the free area is
\begin{equation}\label{eq:volume}
             A=\covol\Lambda-\sum_{i=1}^r\area(C_i)>0.
\end{equation}
Initial conditions have speed one and normalized phase measure
$\dd q\dd\varphi/(2\pi A)$. A selected channel is an isolated,
unobstructed common normal between two obstacle copies. Small facing
patches and a sufficiently short onset collar select exactly the three
successive impacts source--opposite--source.

On success record the first and last signed source arclengths $s,t$
from its normal contact. The corresponding sign on the opposite arc
increases in the same transverse direction at the normal orbit.
Only the source arclengths are observed. Failure is a separate outcome.
At absolute time $T$ the subprobability measure is
\[
 \mu_T(B)=\Prob\{\text{the selected three-impact event by }T,
                                      \ (s,t)\in B\}.
\]
It is not conditioned on success. A local record consists of the two
absolute times $T_2>T_1$ and their measures, on a common strictly active
rectangle, modulo one simultaneous reversal $(s,t)\mapsto(-s,-t)$.
The sign choice is the same at both times. Each record retains its own
contact origin, itinerary selection and grouping of its observations.
No Cartesian position, direction, gap or area is supplied.

The phrase \emph{two density functions} will always mean function-valued
data. Their finite counterparts below are two increasingly fine
histograms, with a number of cells tending to infinity. Neither is a
pair of scalar observations. The local pilot only chooses active times;
it need not estimate the onset to vanishing error.

\begin{theorem}[The smooth local inverse]\label{thm:local-main}
For one selected channel, two strictly active endpoint density functions
$f_1,f_2$ recover its action and area by
\begin{equation}\label{eq:ratio-main}
 W=\frac{T_1f_2-T_2f_1}{f_2-f_1},\qquad
 A=\frac{(T_2-T_1)(-W_{st})}{2\pi(f_2-f_1)}.
\end{equation}
The diagonal gradient and Hessian of $W$ determine both smooth boundary
arcs in their relative placement, up to a common Euclidean isometry.
The gap is $W(0,0)/2$. Under the explicit physical $C^3$, active,
twist and time margins in Proposition~\ref{prop:local-stable}, the
inverse is Lipschitz from the two $C^2$ densities to the intrinsic
$C^2$ arcs, gap and area. Analyticity is unnecessary for this local result.
\end{theorem}

\subsection{An unlabelled collection and a covolume test}
Call a body asymmetric if its Euclidean isometry group is trivial.
Call a table shape-separated if its distinct representatives $C_i$
are pairwise noncongruent, allowing reflections in this comparison.
Repeated observations of translates of the same representative are
allowed. These are conditions on the unknown table, not supplied
boundary shapes or cross-experiment correspondences.

An \emph{unlabelled collection} is a finite multiset of local records.
The pairing of the two windows within a record and the local return
orientation are retained. There is no supplied incidence graph, obstacle
label, cross-channel contact identification, relative sign, period basis,
or integer copy label. Equality allows a permutation of records and an
independent coherent reversal of each record. The underlying channels
are required to visit every obstacle orbit and to form a connected
incidence graph. Coverage and connectedness are properties of the
experiment; they are not inferred for an obstacle never visited.

For analytic boundaries each reconstructed arc determines the entire
boundary image. Congruence classes of the images recover the vertices
of the incidence graph. Matching along a spanning tree places one
representative of each class. A remaining edge gives a translation
$d_e$ between its placed target image and the chosen target representative.
These operations and their independence of all local sign choices are
proved in Section~\ref{sec:descent}. Put
\begin{equation}\label{eq:calibrated-volume}
 V=A+\sum_i\area(C_i),\qquad
 D=[d_{e_1}\ d_{e_2}],\qquad n=\frac{|\det D|}{V},
\end{equation}
where $d_{e_1},d_{e_2}$ are any two independent non-tree displacements.
They need not have been marked by integers.

\begin{theorem}[Unlabelled finite-index descent]\label{thm:descent-main}
Consider an analytic asymmetric, shape-separated table and a connected,
covering unlabelled collection with two independent reconstructed cycle
displacements. The collection determines the obstacle images, incidence
graph, tree placement, free area and period covolume. The number $n$
in \eqref{eq:calibrated-volume} is a positive integer. Every compatible
period lattice is in the following explicit finite list:
\begin{equation}\label{eq:hnf-main}
 \Lambda_B=DB^{-1}\Z^2,\qquad
 B=\begin{pmatrix}a&b\\0&d\end{pmatrix},\quad
 ad=n,\quad a,d>0,\quad 0\le b<d.
\end{equation}
Thus there are at most $\sigma_1(n)=\sum_{d\mid n}d$ compatible tables,
up to a common isometry and permutation of obstacle representatives.
Other cycle, separation and channel constraints can only remove candidates.
If $|\det D|=V$, the whole table and its unmarked lattice are unique.
All assertions use the two endpoint density functions per record, not
counts. No global integer marking or incidence identification is an input.
\end{theorem}

The determinant condition is a check on recovered geometric quantities,
not a supplied primitive pair of integer labels. It is stronger than
merely having two independent cycles. The finite-index qualification is
necessary: Proposition~\ref{prop:ambiguity} constructs, for every odd
prime $p$, $p-1$ nonisometric physical tables with identical selected
local records and calibrated index $p$. Theorem~\ref{thm:descent-main}
also has nonempty open primitive classes with any prescribed finite
number of distinct obstacle shapes.

\begin{theorem}[Finite observations on a primitive class]\label{thm:stat-main}
Fix a compact class of the preceding tables with uniform holomorphic
support-function strips, positive curvature, separation and channel
margins, bounded geometry, quantitative asymmetry and separation between
distinct shape classes. Each selected onset has a known bounded bracket;
the number $J$ of records is bounded. Suppose a reconstructed cycle pair
satisfies $|\det D|=V$. No identity of that pair is supplied to the estimator.

Let $0<\beta\le1$ and assume uniform $C^{2,\beta}$ active-density bounds.
A finite pilot followed by two growing endpoint histograms per record
gives, with probability at least $1-\delta$, whole-table error at most
\begin{equation}\label{eq:rate-main}
 C\left(\frac{\log(CJN/\delta)}{N}\right)^{\beta\theta/(2\beta+6)},
                       \qquad 0<\theta<1,
\end{equation}
using at most $CJN$ independent phase preparations, including failures.
The incidence classes and the primitive period structure are stable for
sufficiently large $N$. The exponent $\theta$, constants and the required
sample size depend on the entire prior class. The loss and all margins
are specified in Section~\ref{sec:stability}. This is a conditional upper
bound, not a minimax exponent or a claim about a count-only sensor.
\end{theorem}

The proof passes from histograms to actual arcs, then to analytic images,
and finally to the finite incidence and period data. It does not estimate
an increasing number of contact jets. The new global step uses the area
normalization in \eqref{eq:ratio-main}, which would be lost by conditioning
on successful trajectories. Symmetric tables, repeated congruent obstacle
representatives and nonprimitive collections retain their earlier marked
or finite-candidate results; they are not silently included in the new
unique unlabelled theorem.
